\subsection{A staged implementation with concrete deliverables}

Begin with a small number of products whose outcomes and economics are well measured, for example one frequently transacted product and one first-adoption product. This tests both recurrence and adoption without trying to resolve every commercial definition simultaneously. The choice is a proposed scope assumption, to be made with product owners using data availability and potential value.

The first stage produces a signed-off outcome dictionary, timestamp and availability rules, entity-group mapping, eligibility logic, missing-data audit and a reproducible monthly panel. Reconcile sampled labels with source systems and CRM records. Establish whether exact transaction dates can be recovered, which determines whether original BG/NBD is a valid benchmark. Record known facility maturities separately from model-inferred signals.

The second stage produces the constructed simple-rule comparisons, a regularized logistic forecast and a partially pooled hazard, followed by a bounded boosting comparison. Its deliverables are chronological predictions, calibration estimates, segment results, paired uncertainty, runtime and memory measurements, and a proposed queue at actual capacity. Web enrichment should enter as an ablation: compare the same model and dates with and without the additional fields. A current web crawl cannot serve as historical data unless archived timestamps establish availability.

The third stage produces a prospectively randomized contact protocol, sample-size analysis, assignment log and complete outcome collection. Historical uplift estimates can nominate strata and hypotheses, but cannot replace the trial. After outcomes mature, estimate assignment effects and a restricted allocation policy using separate fitting and evaluation data. If the action itself has no demonstrated value, changing the ranking algorithm is not a sufficient remedy; reconsider the offer, timing or eligible population.

The fourth stage evaluates the proposed allocation against current prioritization at equal capacity and measures its adoption in the CRM workflow. Expand products, horizons or data sources only after the existing comparison supports the additional complexity. DeepAR and latent maturity models enter as optional research candidates when simpler methods leave a specific, measurable deficiency. The product can remain useful with a classical pooled model if that model meets the commercial criteria.

\subsection{A bounded research protocol before any model campaign}

This document does not execute model training, customer scraping or a sales experiment. Before those activities, the team should set the dataset version, financial outcome, baseline, primary criterion, conditional evaluation situations, uncertainty method, compute and attempt budget, and conditions for stopping or repairing a run. A practical comparison uses a ladder: simple rules, a tuned classical model, a plain nonlinear candidate, then any enhanced candidate. The final test remains untouched throughout tuning.

Every material choice is a hypothesis. Lookback lengths should reflect plausible business cycles and be tested on earlier chronological folds. Regularization controls unstable sparse effects and requires sensitivity checks. Tree depth and shrinkage trade interactions against overfit. A neural candidate needs its own capacity and optimizer search, multiple recorded seeds, early stopping based on an earlier validation window and a sufficient budget for meaningful downstream comparison. A convenient architecture from a retail paper is not a bank-specific default. If the available budget cannot support a fair neural comparison, defer that comparison and retain the validated simpler candidate.

Before expanding computation, run a focused check that the labels, feature timing, loss and probability units are correct. A small successful training run establishes execution, not scientific or commercial adequacy. Store exact configurations, data versions, seeds, runtime, CPU/GPU mode, complete predictions and failure records. A localized infrastructure failure warrants repair within the agreed scope; an invalid target or corrupted outcome record requires revising the experiment before continuing.

\subsection{What remains uncertain}

Three years can support cross-sectional learning but gives little evidence about rare economic regimes. Corporate product demand may be driven by acquisitions, tenders or financing plans absent from transaction histories. External announcements may improve those signals, but their timeliness, coverage and entity matching must be measured. Larger feature sets do not resolve an unobserved demand denominator.

Historical contacts may be too selectively recorded to identify effects. Randomization can address assignment confounding within the experimental population, while interference, incomplete follow-up and changing offers still require attention. Value measured over a short horizon can miss later attrition or credit losses. Where long-run margin cannot yet be observed, report the chosen short-run outcome and the unresolved long-run economics explicitly.

The recommendation is therefore to build the measurable bank-opportunity forecast first, collect credible intervention evidence alongside it, and allocate sales capacity using validated incremental value when that evidence becomes available. The mathematics specifies how to do this coherently. The data and prospective comparisons determine whether it succeeds for the bank.
